\section{近代物理}

\subsection{核心积木}

\begin{结论}{阈能：在质心系算不变质量}
$s=(\sum E)^2-(\sum\mathbf pc)^2$ 是不变量。束流 $m_1$ 打静止靶 $m_2$，产生总静质量 $\sum m_f$：
\[
\重点{E_{1,\mathrm{th}}=\frac{(\sum m_f)^2-m_1^2-m_2^2}{2m_2}\,c^2}\quad(\text{束流总能量}).
\]
$pp\to ppp\bar p$：$E_{\mathrm{th}}=7m_pc^2$，动能 $6m_pc^2\approx5.6\,$GeV。对撞机对头碰撞 $\sqrt s=2E$，固定靶 $\sqrt s\approx\sqrt{2m_2c^2E}$，这就是对撞机的优势。
\end{结论}

\begin{结论}{二体衰变的能量与动量}
静止 $M\to m_1+m_2$：
\[
\重点{E_1=\frac{M^2+m_1^2-m_2^2}{2M}c^2},\qquad |\mathbf p|c=\frac{\sqrt{[M^2-(m_1+m_2)^2][M^2-(m_1-m_2)^2]}}{2M}c^2 .
\]
$\pi^+\to\mu^+\nu$：$p_\mu c\approx29.8\,$MeV，$T_\mu\approx4.1\,$MeV（单一谱线）；三体衰变如 $\beta$ 衰变则为连续谱，这正是中微子假说的来源。
\end{结论}

\begin{结论}{Compton 散射}
\[
\重点{\lambda'-\lambda=\frac{h}{m_ec}(1-\cos\theta)},\qquad \lambda_C=\frac{h}{m_ec}=2.426\ \mathrm{pm},\qquad E'=\frac{E}{1+\frac{E}{m_ec^2}(1-\cos\theta)} .
\]
反冲电子最大动能（$\theta=\pi$）$T_{\max}=\frac{2E^2}{m_ec^2+2E}$，即 Compton 边。逆 Compton：极端相对论电子把低能光子能量提升约 $\frac43\gamma^2$ 倍（各向同性平均）。
\end{结论}

\begin{结论}{相对论 Doppler 与光行差}
源以 $\beta$ 运动，$\theta$ 为观察者系中速度与源到观察者方向的夹角：
\[
\重点{f=\frac{f_0}{\gamma(1-\beta\cos\theta)}},\qquad \text{正对接近 }f=f_0\sqrt{\frac{1+\beta}{1-\beta}},\qquad \cos\theta=\frac{\cos\theta'+\beta}{1+\beta\cos\theta'} .
\]
$\theta=\pi/2$ 时为横向红移 $f_0/\gamma$。源系各向同性发射的一半光子集中在观察者系半角 $\sim1/\gamma$ 的前向锥内（相对论束流）。
\end{结论}

\begin{结论}{快度与匀加速运动}
\重点{\tanh\eta=\beta}，共线速度合成即快度相加；$E=mc^2\cosh\eta$，$p=mc\sinh\eta$。固有加速度 $g$ 恒定（从静止出发）：
\[
\重点{t=\frac cg\sinh\frac{g\tau}{c}},\qquad x=\frac{c^2}{g}\Bigl(\cosh\frac{g\tau}c-1\Bigr).
\]
距离 $c^2/g$ 之后的光信号永远追不上（Rindler 视界）；$g=9.8\,\mathrm{m/s^2}$ 时 $c^2/g\approx0.97$ 光年。
\end{结论}

\begin{结论}{类氢体系的约化质量标度}
\[
\重点{E_n=-\frac{\mu}{m_e}\frac{Z^2}{n^2}\,13.6\ \mathrm{eV}},\qquad r_n=\frac{m_e}{\mu}\frac{n^2}{Z}a_0 .
\]
电子偶素 $\mu=m_e/2$：基态 $-6.8\,$eV；$\mu$ 子氢 $\mu\approx186m_e$：半径缩小约 186 倍。Moseley：$K_\alpha$ 能量 $\approx\frac34(Z-1)^2\times13.6\,$eV（$Z-1$ 为另一个 $1s$ 电子的屏蔽）。
\end{结论}

\begin{结论}{氢原子精细结构与超精细}
\[
\重点{E_{nj}=-\frac{13.6\ \mathrm{eV}}{n^2}\Bigl[1+\frac{\alpha^2}{n^2}\Bigl(\frac{n}{j+\frac12}-\frac34\Bigr)\Bigr]},\qquad \alpha\approx\frac1{137}.
\]
只依赖 $n,j$：$2s_{1/2}$ 与 $2p_{1/2}$ 简并，被 Lamb 位移（QED，约 $1057\,$MHz）打破。基态超精细分裂 $5.9\,\mu$eV，即 $21\,$cm（$1420\,$MHz）谱线；$\Delta E_{\mathrm{hf}}=\frac83g_p\frac{m_e}{m_p}\alpha^2\,\mathrm{Ry}$，$g_p\approx5.59$。
\end{结论}

\begin{结论}{常用数值组合}
\[
\重点{hc=1240\ \mathrm{eV\,nm}},\quad \hbar c=197.3\ \mathrm{MeV\,fm},\quad e^2/(4\pi\varepsilon_0)=1.44\ \mathrm{MeV\,fm},\quad m_ec^2=0.511\ \mathrm{MeV}.
\]
非相对论电子经电压 $V$ 加速：\重点{\lambda=\frac{1.226\ \mathrm{nm}}{\sqrt{V/\mathrm{V}}}}；热中子（$kT\approx25\,$meV）$\lambda\approx0.18\,$nm，适合晶体衍射。$\hbar=6.58\times10^{-22}\,$MeV\,s 把宽度换寿命：$\tau=\hbar/\Gamma$。
\end{结论}

\begin{结论}{Rutherford 散射}
入射动能 $E_k$，电荷 $ze$ 打 $Ze$，最近距离（正碰）$d=\frac{zZe^2}{4\pi\varepsilon_0E_k}$：
\[
\重点{b=\frac d2\cot\frac\theta2},\qquad \frac{\dd\sigma}{\dd\Omega}=\Bigl(\frac{d}{4}\Bigr)^2\frac{1}{\sin^4(\theta/2)} .
\]
$\sigma(\theta>\theta_0)=\pi b(\theta_0)^2$。厚 $x$、数密度 $n$ 的靶：散射比例 $n\,x\,\sigma$（薄靶）；衰减 $I=I_0\ee^{-n\sigma x}$，平均自由程 $1/(n\sigma)$。
\end{结论}

\begin{结论}{半经验质量公式与 $\beta$ 稳定线}
\[
B=a_VA-a_SA^{2/3}-a_C\frac{Z(Z-1)}{A^{1/3}}-a_A\frac{(A-2Z)^2}{A}+\delta,
\]
典型 $a_V\approx15.8$，$a_S\approx18.3$，$a_C\approx0.714$，$a_A\approx23.2\,$MeV。固定 $A$ 对 $Z$ 求原子质量极小（计入 $n$--$p$ 质量差）得 \重点{Z_{\mathrm{stable}}\approx\frac{A}{1.98+0.0155A^{2/3}}}；只用 $B$ 的近似为 $\frac{A}{2+(a_C/2a_A)A^{2/3}}$。$B/A$ 在 $^{56}$Fe 附近最大（约 $8.8\,$MeV），轻核聚变、重核裂变放能；核半径 $R\approx1.2A^{1/3}\,$fm。
\end{结论}

\begin{结论}{衰变链（Bateman 两步）}
$1\to2\to3$，$N_2(0)=0$：
\[
\重点{N_2(t)=\frac{\lambda_1N_1(0)}{\lambda_2-\lambda_1}\bigl(\ee^{-\lambda_1t}-\ee^{-\lambda_2t}\bigr)},\qquad t_{\max}=\frac{\ln(\lambda_2/\lambda_1)}{\lambda_2-\lambda_1}.
\]
$\lambda_1\ll\lambda_2$：长期平衡，活度 $\lambda_1N_1=\lambda_2N_2$；$\lambda_1<\lambda_2$ 但同量级：暂时平衡，$\frac{A_2}{A_1}\to\frac{\lambda_2}{\lambda_2-\lambda_1}$。
\end{结论}

\begin{结论}{电偶极跃迁选择定则}
\[
\重点{\Delta l=\pm1,\quad \Delta m_l=0,\pm1,\quad \Delta S=0,\quad \text{宇称改变}};\qquad \Delta J=0,\pm1\ (J=0\not\to0).
\]
速率 $\Gamma=\frac{\omega^3|\langle f|e\mathbf r|i\rangle|^2}{3\pi\varepsilon_0\hbar c^3}$；氢 $2p\to1s$ 寿命 $1.6\,$ns，$2s$ 只能双光子衰变（约 $0.12\,$s）。$\Delta m=0$ 为 $\pi$ 光（沿场线偏振），$\pm1$ 为 $\sigma^\pm$。
\end{结论}

\begin{结论}{弱场广义相对论的四个数}
\[
\重点{\frac{\Delta f}{f}=\frac{gh}{c^2}},\quad \delta\phi_{\text{光偏折}}=\frac{4GM}{c^2b},\quad \Delta\varpi=\frac{6\pi GM}{c^2a(1-e^2)}\ (\text{每圈}),\quad r_s=\frac{2GM}{c^2}.
\]
太阳边缘偏折 $1.75''$，水星进动每世纪 $43''$。GPS 卫星钟：狭义相对论 $-7\,\mu$s/天，引力 $+45\,\mu$s/天，净快约 $38\,\mu$s/天。
\end{结论}

\参考{French, \emph{Special Relativity}；Taylor--Wheeler, \emph{Spacetime Physics}（MIT 8.033 / 8.20）；Krane, \emph{Introductory Nuclear Physics}；Griffiths, \emph{Introduction to Elementary Particles}；Foot, \emph{Atomic Physics}；PDG 常数表。}
